\documentclass[aspectratio=169,10pt]{beamer}
\usepackage[T1]{fontenc}
\usepackage{lmodern}
\usepackage{booktabs}
\definecolor{courseblue}{RGB}{27,57,85}
\setbeamercolor{structure}{fg=courseblue}
\setbeamercolor{frametitle}{fg=courseblue}
\setbeamerfont{frametitle}{size=\large,series=\bfseries}
\setbeamertemplate{navigation symbols}{}
\setbeamertemplate{footline}{\hspace{0.35cm}\scriptsize IIT Gandhinagar\hfill CE 637 Course Project\hspace{0.35cm}\vspace{0.15cm}}
\setbeamersize{text margin left=0.55cm,text margin right=0.55cm}
\begin{document}
\begin{frame}[t]{Food Grain Distribution in Gujarat: Summary}
{\scriptsize Abhitej Singh Bhullar (23110009) \quad Radhika Purohit (24110280) \hfill November 2026}\par\vspace{0.10cm}\small
\begin{columns}[T,onlytextwidth]
\begin{column}{0.49\textwidth}\small
\textbf{The decision}\par
Allocate wheat and rice from four depots to six districts at minimum cost, within stock and loading limits.
\vspace{0.15cm}

\textbf{Model and method}
\begin{itemize}\setlength{\itemsep}{1pt}
\item 48 shipment variables; 24 constraints.
\item Two-Phase Simplex: 56 pivots; same cost as SciPy HiGHS.
\end{itemize}
\vspace{0.1cm}
\textbf{Why Surat needs four depots}\par
{\small
\begin{tabular}{lr}
\toprule
Supply depot & Surat allocation (MT)\\
\midrule
Sabarmati & 11,000\\
Chhani & 7,000\\
Ghanteshwar & 12,000\\
Gandhidham & 4,000\\
\bottomrule
\end{tabular}}
\end{column}
\begin{column}{0.47\textwidth}\small
\textbf{Main results}\par\vspace{0.1cm}
\begin{tabular}{@{}lr@{}}
Demand met & \textbf{132,000 MT}\\[2pt]
Monthly cost & \textbf{Rs.\ 5.6983 crore}\\[2pt]
Average cost & \textbf{Rs.\ 431.69/MT}\\[2pt]
Active routes & \textbf{9 of 24}\\
\end{tabular}
\vspace{0.25cm}

\textbf{Chhani loading-capacity expansion}
\begin{itemize}\setlength{\itemsep}{2pt}
\item Raise monthly capacity from 28,000 to 34,000 MT.
\item Save \textbf{Rs.\ 55.05 lakh/month}, or \textbf{Rs.\ 6.61 crore/year}.
\item Baseline shadow price: Rs.\ 1,022.90/MT; no benefit beyond 34,000 MT stock.
\end{itemize}
\textbf{Recommendation:} Compare loading-capacity expansion costs with the estimated savings.
\end{column}
\end{columns}
\vspace{0.05cm}
{\scriptsize Inputs: assumed stocks, demand, loading capacity and freight rates; OSRM road distances.}
\end{frame}
\end{document}
